% 16-boundaries.tex: what Polaris deliberately refuses to become.
\section{What Polaris refuses to become}
\label{sec:boundaries}

Powerful identity infrastructure carries risks that no amount of correctness removes, because they
are risks of what a correct system can be used for. The repository names these as design and
governance constraints rather than as politics, and this section states them the same way. Each is
a way a working Polaris could be turned into something its constitution forbids.

\begin{description}
  \item[Identity must not become money.] A credential that carries spending authority inherits
    programmability, and every constraint anyone wants placed on spending accretes onto the identity
    token, until the system can be told whom to refuse at a fuel pump. C10 forbids a monetary claim
    in the schema; if a value system is ever built, it lives in a separate database that consumes
    verification proofs over a boundary, and the boundary is what is load-bearing. This is the single
    most consequential architectural decision in the repository, and it is the one that never
    expires with any phase.
  \item[Privacy is not merely secrecy.] A system can know very little about a person and still be
    coercive, if every doorway requires its authorization. The zero-knowledge default bounds what is
    recorded; it does not bound what a society demands. The tiered-enrolment layer exists for this:
    exemption is a first-class, one-row state, and the not-enrolled cannot be enumerated by an
    accident of schema. The repository's own threat model names the strongest attack on that layer as
    making non-enrolment civically impossible from outside, no token, no healthcare, and says plainly
    that the schema cannot stop it, only make it a named act.
  \item[Technical federation is not social decentralisation.] Polaris is federated in code: each
    authority its own key, no central store, no transitive trust. Two instances run by one author on
    one machine demonstrate the protocol and nothing about who runs it. Whether authorities, witnesses
    and relying parties are in fact independent is a property of a deployment, not of the repository,
    and this paper marks every such claim as requiring external validation.
  \item[Revocation must not become civil erasure.] A revoked token is a terminal state in an
    append-only record, not a deletion; lost tokens stay lost, and their data is not erased.
    Succession is by reference. Mass revocation is bounded by a ceiling and a co-signer. Erasure, when
    a jurisdiction requires it, is pseudonymisation with an append-only record of who, when and why,
    never of the prior name. An authority can end a person's credential; the schema refuses to let it
    end the person's history.
  \item[Relying-party correlation is different from issuer unlinkability.] The issuer's database
    cannot rebuild the zero-knowledge verification graph. Two relying parties can join their logs by
    the credential value. Both sentences are true at once, and the second is a permanent, documented
    property (Section~\ref{sec:privacy}). A reader who hears \emph{unlinkable} and imagines the
    second has been misled, which is why the repository states it positively.
  \item[No population graph.] The Atlas is bounded and never locates a zero-knowledge event; Athena
    is forbidden a person node, a stable subject surrogate and a subject traversal; the schema carries
    no attribute to filter a population by; quotas and alerts are keyed by agency, never by person.
    The refusal is structural because a population graph, once built, outlives the intentions of
    whoever built it.
  \item[Machine learning must not turn identity infrastructure into scoring.] The proposed learning
    layer may learn how the identity system behaves and must never learn how people behave. Fraud
    or risk scores for individuals, behavioural profiles, relationship inference, eligibility or
    likelihood-of-revocation scoring, person-level clustering or embeddings, and model-driven denial
    of rights are presumptively prohibited, and a recommendation is evidence, never authority.
\end{description}

The repository also lists permanent non-goals that do not expire with any roadmap phase: no payment
rails, no central biometric database, no population-scale attribute filtering or analytics, no social
scoring, no commercial or advertising use of verification data, and no real personal data before a
consent framework exists. Every one of them is a way the system could be made more useful to someone,
and every one is refused on sight.

\analogy{The laws above the town bind the town. The hall may not mint coin. The register may not
be turned into a list of who is absent. The gates may not be closed on a neighbour on the word of a
neighbour's neighbour. A name struck from the register stays legible as struck. And no clerk, however
clever, may be set to guessing which townspeople are likely to cause trouble.}
